\subsection{An exact executable classifier for commensurate energies}\label{sec:exactclassifier}
The general thermal test involves transcendental quantities. An important
subfamily admits a terminating exact implementation, including tangencies,
endpoint roots, and algebraic slow eigenvectors. Suppose
\[
E_i=m_i\Delta,\qquad m_1<m_2<m_3<m_4\ \hbox{integers},\qquad
b=e^{-\beta_b\Delta}\in\mathbb Q\cap(0,1),
\]
and the original transition rates are rational. A shift makes $m_1=0$.
Put $x=e^{-\beta\Delta}$; the positive-temperature range is $0<x<1$,
the bath is $x=b$, and $Z(x)=\sum_i x^{m_i}$. Rational rates and bath
weights make the generator rational, and its slow eigenvectors algebraic.

For a simple slow mode define $N(x)=\sum_iw_i x^{m_i}$. Then
\[
b_2(x)=\frac{N(x)}{Z(x)},\qquad
\operatorname{sgn}\frac{db_2}{dx}
=\operatorname{sgn}\{N'Z-NZ'\}.
\]
All coefficients belong to the real algebraic field generated by the slow
eigenvalue. For a degenerate slow space use the amplitude form of
Section~\ref{sec:degenerateextension}, so $Q=N/Z^2$ and the derivative
numerator is $N'Z-2NZ'$. The projector is rational in this case: a repeated
root of the rational nonzero cubic characteristic polynomial must be
rational, and the remaining eigenvalue is rational too. Polynomial spectral
projection therefore avoids any square roots of the bath weights.

Factor the derivative numerator into square-free factors, retaining those
of odd multiplicity. A Sturm count of their roots in $(b,1)$ and $(0,b)$
decides ordinary cooling and heating, respectively: an effect exists exactly
when the corresponding count is nonzero. Indeed, odd-multiplicity roots are
precisely derivative sign changes. For the scalar overlap, a fixed derivative
sign and the bath zero imply normal amplitude ordering; conversely normal
absolute-amplitude ordering forces a fixed sign of the overlap and its
derivative. For the nonnegative degenerate amplitude, the derivative has
the normal sign near its isolated bath zero, so any reversal requires and
is implied by an interior sign change. Even-multiplicity stationary points
alone do not count as effects.

A separate square-free Sturm count of $N$ gives the numbers of strong
preparations on the two intervals. Factors at $x=0,b,1$ are explicitly
removed before interval counts: infinite and zero temperature endpoints,
and the bath itself, are not finite nonbath preparations.

\texttt{exact\_commensurate\_classifier.py} implements these operations for
all three possible slow multiplicities and permits zero rates when the
transition graph remains connected. Signs in an algebraic coefficient field
are determined by rational isolating intervals for its specified real
embedding, refined until interval evaluation excludes zero. Exact field
reduction detects zero first, so refinement terminates for every nonzero
coefficient. This avoids treating a polynomial representation's leading sign
as the sign of its algebraic value. Sturm and square-free arithmetic are
classical algorithms; the contribution here is a reproducible implementation
of the stated physical classification with all equality cases retained.

Sixteen checks cover irrational and scaled cubic slow eigenvalues, double and triple
slow degeneracy, strong roots without sign changes, stationary tangencies
without an effect, bath-threshold equality, excluded temperature endpoints,
nonuniform commensurate spacing, and connected generators with zero rates.
An explicit sign-regression test ensures that the real root $\alpha$ of
$\alpha^3-\alpha-1=0$ satisfies $\alpha-2<0$. The implementation uses exact
SymPy polynomial arithmetic and rational interval refinement, not temperature
sampling or a floating-point eigenvector. Its hypotheses on rationality and
commensurability are substantive; they are not assumed in the preceding
analytic classification.
